%!TEX root = ../main.tex

\subsection{Соотношения баланса энергии}
\label{ssec:energy_balance}

Будем рассматривать следующую постановку задачи. Пусть образец вещества имеет форму параллелепипеда, то есть область решения $\clOmega = [0, L]_x \times [0, H]_y \times [0, D]_z$. Образец зажат между двумя проводящими пластинами $y = 0$ и $y = H$ под напряжением.

\begin{subequations}
\label{cond:boundary_basic}
На <<боковых>> гранях $\partial \Omega$ зададим однородное условие Неймана на электрический потенциал:
\begin{equation}
    \partvn{\Phi} \bigg|_{ \{ x = 0, L \} \cup \{ z = 0, D \} } = 0;
\end{equation}
а также условие Дирихле либо Неймана на фазовое поле:
\begin{equation}
    \phi|_{ \{ x = 0, L \} \cup \{ z = 0, D \} } = 1 \quad \text{либо} \quad \partvn{\phi} \bigg|_{ \{ x = 0, L \} \cup \{ z = 0, D \} } = 0.
\end{equation}
На верхней и нижней грани считаем заданным однородное условие Неймана на фазовое поле:
\begin{equation}
    \partvn{\phi} \bigg|_{y = 0, H} = 0.
\end{equation}
Так как пластины из идеального проводника, то
\begin{equation}
    \Phi|_{y = 0} = \Phi^+(t), \quad \Phi|_{y = H} = \Phi^-(t),
\end{equation}
где $\Phi^+(t)$ и $\Phi^-(t)$~-- функции времени. Потенциал на обкладках может меняться как из-за притока заряда извне, так и за счет изменения сопротивления образца между ними.
\end{subequations}

На обкладках сконцентрированы следующие поверхностные заряды:
\begin{equation}
    q^-(t) = \int\limits_{y = 0} \epsilon \cdot (\vE, \vn) dS = \int\limits_{y = 0} \epsilon \party{\Phi} dS, \qquad q^+(t) = \int\limits_{y = H} \epsilon \cdot (\vE, \vn) dS = -\int\limits_{y = H} \epsilon \party{\Phi} dS.
    \label{eq:charge}
\end{equation}
Интегрируя уравнение \eqref{eq:electrostatic_c}, по формуле Гаусса--Остроградского
\[
    \intOmega \rho d \vx = \intOmega \Div(-\epsilon \nabla \Phi) = \int\limits_{y = 0} \epsilon \party{\Phi} dS - \int\limits_{y = H} \epsilon \party{\Phi} dS = q^- + q^+.
\]

Можно считать $\Phi^- = 0$, то есть уровнем нулевого потенциала. Заданное $\Phi^+$ замыкает граничные условия, тогда однозначно определяется $\Phi$, а значит, и $q^-$, $q^+$ (по формуле \eqref{eq:charge}). Верно и обратное: заданное $q^-$ определяет $\Phi^+$. Для этого рассмотрим однородную ($\rho \equiv 0$) задачу с $\Phi^+ = 1$ и неоднородную с исходным $\rho$ и $\Phi^+ = 0$. Линейная комбинация их решений с коэффициентами $k$, $1$ даст исходное значение $\rho$ и произвольное~-- $q^-$.

\begin{proposition}
\label{prop:minimization}
Уравнение \eqref{eq:phi} динамики фазового поля удовлетворяет принципу минимизации функционала внутренней энергии системы
\begin{equation}
    W = \intOmega w d \vx
    \label{eq:energy_inner}
\end{equation}
на касательном пространстве, определяемом системой уравнений \eqref{eq:electrostatic}, при мгновенном постоянстве $q^-$.
\end{proposition}
\begin{proof}
Переходя к мгновенным вариациям в уравнениях \eqref{eq:electrostatic_c} и \eqref{eq:charge} (с учетом постоянства $q^\pm$), получаем следующие уравнения касательного пространства:
\begin{gather}
    \Div \left( \delta \epsilon \nabla \Phi + \epsilon \nabla \delta \Phi \right) = 0,
    \label{eq:tangent_instant_rho} \\
    \int\limits_{y = 0, H} \left( \delta \epsilon \party{\Phi} + \epsilon \party{\delta \Phi} \right) dS = 0.
    \label{eq:tangent_instant_charge}
\end{gather}

Рассмотрим вариацию функционала $W$:
\begin{equation}
    \delta W = \delta W_E + \intOmega \left[ \cfrac{\Gamma}{l^2} f'(\phi) \delta \phi + \cfrac{\Gamma}{2} (\nabla \delta \phi, \nabla \phi) + \beta \Gamma l^2 \left( \nabla \delta \phi, \| \nabla \phi \|_2^2 \nabla \phi \right) \right] d \vx.
    \label{eq:delta_energy}
\end{equation}
Используя известное дифференциальное тождество
\[
    \Div(\xi \nabla \eta) = (\nabla \xi, \nabla \eta) + \xi \Delta \eta,
\]
а также учитывая поставленные граничные условия типа Дирихле либо однородные типа Неймана, для слагаемых с дифференциальными операторами стандартным образом получаем
\begin{equation}
    \intOmega \left[ \cfrac{\Gamma}{2} (\nabla \delta \phi, \nabla \phi) + \beta \Gamma l^2 \left( \nabla \delta \phi, \| \nabla \phi \|_2^2 \nabla \phi \right) \right] d \vx = -\intOmega \delta \phi \left[ \cfrac{\Gamma}{2} \Delta \phi + \beta \Gamma l^2 \Div \left(\| \nabla \phi \|_2^2 \nabla \phi \right) \right] d \vx.
    \label{tmp:delta_energy_differential}
\end{equation}

Осталось преобразовать $\delta W_E$.
\[
    \delta W_E = \intOmega \left[ \epsilon \cdot (\nabla \Phi, \nabla \delta \Phi) + \half \delta \epsilon \cdot (\nabla \Phi, \nabla \Phi) \right] d \vx,
\]
где для первого слагаемого
\begin{equation}
    \intOmega \epsilon \cdot (\nabla \Phi, \nabla \delta \Phi) d \vx = \intOmega \Div (\Phi \epsilon \nabla \delta \Phi) d \vx - \intOmega \Phi \Div (\epsilon \nabla \delta \Phi) d \vx.
    \label{tmp:detla_energy_electrical_first}
\end{equation}
Далее, используя формулу Гаусса--Остроградского, учитывая равенство \eqref{eq:tangent_instant_charge} и $\Phi^- = 0$,
\begin{multline*}
    \intOmega \Div (\Phi \epsilon \nabla \delta \Phi) d \vx = \intpartOmega \Phi \epsilon \cdot (\nabla \delta \Phi, \vn) dS = -\Phi^+ \int\limits_{y = 0} \epsilon \party{\delta \Phi} dS = \Phi^+ \int\limits_{y = 0} \delta \epsilon \party{\Phi} dS = \\ = -\intpartOmega \Phi \delta \epsilon \cdot (\nabla \Phi, \vn) = -\intOmega \Div(\Phi \delta \epsilon \nabla \Phi) d \vx;
\end{multline*}
а также, используя равенство \eqref{eq:tangent_instant_rho},
\[
    \intOmega \Phi \Div (\epsilon \nabla \delta \Phi) d \vx = -\intOmega \Phi \Div(\delta \epsilon \nabla \Phi).
\]
Возвращаясь к выражению \eqref{tmp:detla_energy_electrical_first}, имеем
\[
    \intOmega \epsilon \cdot (\nabla \Phi, \nabla \delta \Phi) d \vx = -\intOmega \left[ \Div(\Phi \delta \epsilon \nabla \Phi) - \Phi \Div(\delta \epsilon \nabla \Phi) \right] d \vx = -\intOmega \delta \epsilon \cdot (\nabla \Phi, \nabla \Phi) d \vx.
\]
Таким образом,
\begin{equation}
    \delta W_E = -\intOmega \delta \epsilon \cdot (\nabla \Phi, \nabla \Phi) d \vx + \half \intOmega \delta \epsilon \cdot (\nabla \Phi, \nabla \Phi) d \vx = -\half \intOmega \delta \epsilon \cdot (\nabla \Phi, \nabla \Phi) d \vx.
    \label{tmp:delta_energy_electrical}
\end{equation}

Наконец, подставляя \eqref{tmp:delta_energy_differential} и \eqref{tmp:delta_energy_electrical} в \eqref{eq:delta_energy}, получаем
\[
    \cfrac{\delta W}{\delta \phi} = -\left[ \half \epsilon \cdot (\nabla \Phi, \nabla \Phi) + \cfrac{\Gamma}{l^2} f' + \cfrac{\Gamma}{2} \Delta \phi + \beta \Gamma l^2 \Div \left(\| \nabla \phi \|_2^2 \nabla \phi \right) \right],
\]
где в квадратных скобках стоит правая часть уравнения \eqref{eq:phi} динамики фазового поля.
\end{proof}

\begin{proposition}
\label{prop:energy_electrical_balance}
Справедливо следующее уравнение баланса энергии $W_E$ электрического поля:
\begin{equation}
    \cfrac{d W_E}{dt} = -\half \intOmega \partt{\epsilon} (\nabla \Phi, \nabla \Phi) d \vx - \Phi^+ \left( \partt{q^-} + \int\limits_{y = 0} \sigma \party{\Phi} dS \right) - \intOmega Q_j d \vx,
    \label{eq:energy_electrical_balance}
\end{equation}
где $Q_j$, определяемое формулой \eqref{eq:heat}, есть тепло, выделяемое при протекании тока в среде.
\end{proposition}
\begin{proof}
Повторяется схема доказательства утверждения \ref{prop:minimization}.

Продифференцируем по времени уравнения \eqref{eq:electrostatic_c} и \eqref{eq:charge}, получим
\begin{gather}
    \Div \left( \partt{\epsilon} \nabla \Phi + \epsilon \nabla \partt{\Phi} \right) = -\partt{\rho},
    \label{eq:tangent_temporal_rho} \\
    \int\limits_{y = 0} \left( \partt{\epsilon} \party{\Phi} + \epsilon \partt{} \party{\Phi} \right) dS = \partt{q^-}.
    \label{eq:tangent_temporal_charge}
\end{gather}

Продифференцируем по времени функционал $W_E$:
\begin{equation}
    \cfrac{d W_E}{dt} = \intOmega \left[ \epsilon \cdot \left( \nabla \Phi, \nabla \partt{\Phi} \right) + \half \partt{\epsilon} (\nabla \Phi, \nabla \Phi) \right] d \vx.
    \label{tmp:energy_electrical_deriv}
\end{equation}
Преобразуя первое слагаемое с учетом \eqref{eq:tangent_temporal_rho}, \eqref{eq:tangent_temporal_charge} тем же образом, что для $\partial W_E$ в доказательстве утверждения \ref{prop:minimization}, имеем
\begin{equation}
    \intOmega \epsilon \cdot \left( \nabla \Phi, \nabla \partt{\Phi} \right) d \vx = -\intOmega \partt{\epsilon} (\nabla \Phi, \nabla \Phi) d \vx - \Phi^+ \partt{q^-} + \intOmega \Phi \partt{\rho} d \vx.
    \label{tmp:energy_electrical_deriv_first}
\end{equation}
Выразим производную заряда из уравнения \eqref{eq:electrostatic_b} и подставим в последнее слагаемое выражения выше:
\begin{multline*}
    \intOmega \Phi \partt{\rho} d \vx = \intOmega \Phi \Div (\sigma \nabla \Phi) d \vx = \intpartOmega \Phi \sigma \cdot (\nabla \Phi, \vn) dS - \intOmega \sigma \cdot (\nabla \Phi, \nabla \Phi) d \vx = \\ = -\Phi^+ \int\limits_{y = 0} \sigma \party{\Phi} dS - \intOmega Q_j d \vx.
\end{multline*}

В итоге, подставляя два последних равенства выше в выражение \eqref{tmp:energy_electrical_deriv}, получаем уравнение \eqref{eq:energy_electrical_balance} баланса электрической энергии.
\end{proof}

\begin{proposition}
\label{prop:energy_inner_balance}
Справедливо следующее уравнение баланса внутренней энергии $W$ системы:
\begin{equation}
    \cfrac{dW}{dt} = -\cfrac{1}{m} \intOmega \left( \partt{\phi} \right)^2 d \vx - \Phi^+ \left( \partt{q^-} + \int\limits_{y = 0} \sigma \party{\Phi} dS \right) - \intOmega Q_j d \vx.
    \label{eq:energy_inner_balance}
\end{equation}
\end{proposition}
\begin{proof}
Так как плотность $w$ внутренней энергии получается прибавлением к $w_E$ $\Gamma / l^2 \cdot [1 - f(\phi)]$ и еще двух дифференциальных слагаемых (см. \eqref{eq:energy_density_inner}), то проделав при дифференцировании этих слагаемых те же стандартные действия, что в начале доказательства \ref{prop:minimization}, получим
\begin{multline*}
    \cfrac{d}{dt} \intOmega \left[ \cfrac{\Gamma}{l^2}[1 - f(\phi)] + \cfrac{\Gamma}{4} \| \nabla \phi \|_2^2 + \cfrac{\beta \Gamma l^2}{4} \| \nabla \phi \|_2^4 \right] d \vx = \\ = - \intOmega \partt{\phi} \left[ \cfrac{\Gamma}{l^2} f' + \cfrac{\Gamma}{2} \Delta \phi + \beta \Gamma l^2 \Div \left( \| \nabla \phi \|_2^2 \nabla \phi \right) \right] d \vx.
\end{multline*}
Прибавляя полученную производную к равенству \eqref{eq:energy_electrical_balance} и выделяя под интегралом правую часть уравнения \eqref{eq:phi} динамики фазового поля, получаем уравнение \eqref{eq:energy_inner_balance}.
\end{proof}

\begin{remark}
\label{rem:charge_conservative}
Если считать, что приток заряда на обкладки извне отсутствует, то
\[
    \cfrac{d q^-}{dt} = -\int\limits_{y = 0} \sigma \party{\Phi} dS.
\]
В этом случае уравнения \eqref{eq:energy_electrical_balance} и \eqref{eq:energy_inner_balance} соответственно принимают вид
\begin{align*}
    \cfrac{d W_E}{dt} & = -\half \intOmega \partt{\epsilon} (\nabla \Phi, \nabla \Phi) d \vx - \intOmega Q_j d \vx, \\
    \cfrac{dW}{dt} & = -\cfrac{1}{m} \intOmega \left( \partt{\phi} \right)^2 d \vx - \intOmega Q_j d \vx.
\end{align*}
Легко видеть, что в таком случае $dW / dt \leq 0$~-- $W$ не возрастает.
\end{remark}

\begin{remark}
\label{rem:charge_constant}
Если верно $\sigma = k\epsilon$, $k$~-- константа, то из уравнений \eqref{eq:electrostatic_b}, \eqref{eq:electrostatic_c} следует
\[
    \partt{\rho} + k \rho = 0, \quad \rho(\vx, t) = \rho(\vx, 0) e^{-kt}.
\]
Если притом начальная плотность заряда $\rho(\vx, 0) = \rho_0(\vx) = 0$ (и, следовательно, $\rho \equiv 0$), а также заряды на обкладках поддерживаются постоянными: $q^-(t) = -q^+(t) = \Const$, то балансовые соотношения \eqref{eq:energy_electrical_balance} и \eqref{eq:energy_inner_balance} соответственно принимают вид
\begin{align*}
    \cfrac{d W_E}{dt} & = -\half \intOmega \partt{\epsilon} (\nabla \Phi, \nabla \Phi) d \vx, \\
    \cfrac{dW}{dt} & = -\cfrac{1}{m} \intOmega \left( \partt{\phi} \right)^2 d \vx,
\end{align*}
так как второе и третье слагаемые в соотношении \eqref{tmp:energy_electrical_deriv_first} становятся нулевыми. Очевидным образом, $W$ не возрастает.

Более того, в этом случае $W_E$ и $W$ ведут себя подобно консервативным величинам. Действительно, в такой постановке задачи потенциал $\Phi$ однозначно определяется фазовым полем $\phi$, а значит, справедливы функциональные зависимости $W_E = W_E[\phi]$, $W = W[\phi]$. В строго физическом смысле, естественно, $W_E$ и $W$ не консервативны, так как происходит потеря энергии в виде тепла; однако она <<компенсируется>> за счет нагнетания заряда на обкладки.
\end{remark}